% -*- TeX-master: "../main.tex" -*-
\chapter{Projekt Interfejsu}
Prace nad interfejsem rozpoczęliśmy od przygotowania wstępnego szkicu aplikacji \ref{fig:init-proj}, wykorzystując do tego narzędzie Discord Whiteboard.
\newline\newline
\MyFigure{images/initial.png}{Początkowy projekt}{fig:init-proj}{Discord Whiteboard}
Dzięki temu szkicowi dokładnie wiedzieliśmy, jakie ekrany będzie zawierać nasza aplikacja oraz jak będą wyglądały przejścia między nimi. Kluczową kwestią była dla nas spójność interfejsu. Kierowaliśmy się przy tym zasadą: za każdą konkretną czynność powinien odpowiadać wyłącznie jeden dedykowany element na ekranie.
\section{Pop-up}
Do zaprezentowania użytkownikowi predefiniowanych opcji wykorzystaliśmy okna typu pop-up. W naszej aplikacji służą one między innymi do wyboru metody importu danych (rysunek \ref{fig:import_screen}) czy określenia rodzaju konta (rysunek \ref{fig:settings_screen_acc}).
\MyFigure[0.3\textwidth]{images/elementy/popup.png}{Przykład Pop-up'u w naszej aplikacji.}{fig:popup}{\parencite{makieta}}
\section{Jumper}
Podczas projektowania interfejsu często wykorzystywaliśmy szuflady wysuwające się z dołu ekranu (nazywanych w naszej makiecie jumperami \ref{fig:jumper}). Jest to alternatywa dla znanych chociażby z systemu Android okien dialogowych, która - w naszej opinii - jest bardziej przyjazna dla użytkownika.
\MyFigure[0.3\textwidth]{images/elementy/jumper.png}{Przykład jumper'a w naszej aplikacji.}{fig:jumper}{\parencite{makieta}}
Ten element UI wysuwa się od spodu, delikatnie przyciemniając tło. Wykorzystywany jest przez nas głównie do pobierania dodatkowych danych od użytkownika (rysunek \ref{fig:new_proj_screen}) oraz wyświetlania informacji, którymi jest zainteresowany w danej chwili, takich jak na przykład dokumentacja (rysunek \ref{fig:code_doc_screen}).
\newline\newline
Ogromną zaletą tego rozwiązania jest intuicyjny sposób zamykania. Wystarczy, że użytkownik przesunie palcem w dół lub dotknie przyciemnionego obszaru poza panelem. Jest to dobry sposób na uniknięcie pomyłkowego kliknięcia. Co więcej, panel nie zasłania środkowej części ekranu, a jedynie jego dolny fragment. Dzięki temu użytkownik, czytając dokumentację, nie traci z oczu kodu, który właśnie pisał.
\newline\newline
Rozwiązanie to jest również wygodniejsze dla osób obsługujących aplikację jedną ręką, ponieważ pozwala na zamknięcie okna bez konieczności niewygodnego sięgania palcem do przycisku "anuluj". W efekcie nasza aplikacja jest nie tylko przystępniejsza dla szerszego grona odbiorców, ale również jej prezentacja jest o wiele nowocześniejsza niż konkurencyjne rozwiązania, takie jak na przykład Carnets.
\section{Slider}
Ostatnim istotnym elementem interfejsu naszej aplikacji, o działaniu zbliżonym do jumpera, jest tzw. slider \ref{fig:slider}. Elementy tego typu są często spotykane chociażby w aplikacjach na system Android. Slider funkcjonuje niemal identycznie co jumper, z tą różnicą, że wysuwa się z boku ekranu.
\MyFigure[0.3\textwidth]{images/elementy/slider.png}{Przykład slider'a w naszej aplikacji.}{fig:slider}{\parencite{makieta}}
Wykorzystujemy ten element jako substytut osobnego okna aplikacji. Pozwala on na szybki podgląd informacji o projekcie bez konieczności opuszczania aktualnego widoku, takich jak struktura plików czy błędy w kodzie.
\newline\newline
Dzięki połączeniu tych dwóch elementów (slider \ref{fig:slider} i jumper \ref{fig:jumper}), użytkownik programując cały czas pozostaje w jednym oknie edytora. Opuszcza je wyłącznie wtedy, gdy chce sprawdzić stan swojego projektu.
